\documentclass[12pt,a4paper]{article}

% ================================================
% PACKAGES
% ================================================
\usepackage[utf8]{inputenc}
\usepackage[T1]{fontenc}
\usepackage[turkish]{babel}
\usepackage{geometry}
\usepackage{graphicx}
\usepackage{listings}
\usepackage{xcolor}
\usepackage{hyperref}
\usepackage{titlesec}
\usepackage{booktabs}
\usepackage{longtable}
\usepackage{array}
\usepackage{enumitem}
\usepackage{fancyhdr}
\usepackage{mdframed}
\usepackage{tcolorbox}
\usepackage{fontawesome5}
\usepackage{amsmath}
\usepackage{microtype}

\geometry{margin=2.5cm, top=3cm, bottom=3cm}

% ================================================
% COLORS
% ================================================
\definecolor{indigo}{HTML}{6366F1}
\definecolor{slate900}{HTML}{0F172A}
\definecolor{slate800}{HTML}{1E293B}
\definecolor{slate700}{HTML}{334155}
\definecolor{slate400}{HTML}{94A3B8}
\definecolor{purple}{HTML}{8B5CF6}
\definecolor{codebg}{HTML}{1E293B}
\definecolor{codeframe}{HTML}{334155}
\definecolor{codecomment}{HTML}{64748B}
\definecolor{codestring}{HTML}{A5F3FC}
\definecolor{codekeyword}{HTML}{A78BFA}
\definecolor{codefunction}{HTML}{6EE7B7}
\definecolor{deleted}{HTML}{F87171}
\definecolor{added}{HTML}{6EE7B7}

% ================================================
% CODE LISTINGS STYLE
% ================================================
\lstdefinestyle{saas}{
  backgroundcolor=\color{codebg},
  commentstyle=\color{codecomment}\itshape,
  keywordstyle=\color{codekeyword}\bfseries,
  stringstyle=\color{codestring},
  basicstyle=\ttfamily\footnotesize\color{white},
  breakatwhitespace=false,
  breaklines=true,
  captionpos=b,
  keepspaces=true,
  numbers=left,
  numberstyle=\tiny\color{slate400},
  numbersep=8pt,
  showspaces=false,
  showstringspaces=false,
  showtabs=false,
  tabsize=2,
  frame=single,
  rulecolor=\color{codeframe},
  framesep=4pt,
  xleftmargin=15pt,
  xrightmargin=5pt,
}

\lstdefinestyle{diff}{
  basicstyle=\ttfamily\scriptsize,
  backgroundcolor=\color{codebg},
  frame=single,
  rulecolor=\color{codeframe},
  breaklines=true,
  moredelim=**[is][\color{added}]{@@+}{+@@},
  moredelim=**[is][\color{deleted}]{@@-}{-@@},
}

\lstset{style=saas}
\lstdefinelanguage{JSX}{
  keywords={import,export,default,const,let,var,function,return,if,else,for,while,class,new,this,from,of,in,async,await,true,false,null,undefined},
  morestring=[b]",
  morestring=[b]',
  morestring=[b]`,
  morecomment=[l]{//},
  morecomment=[s]{/*}{*/},
  sensitive=true
}
\lstdefinelanguage{CSS}{
  keywords={root,data-theme,body,hover},
  morestring=[b]",
  morecomment=[s]{/*}{*/},
  sensitive=false
}

% ================================================
% PAGE STYLE
% ================================================
\pagestyle{fancy}
\fancyhf{}
\fancyhead[L]{\small\color{slate400}CodeAlchemist — SaaS Dönüşüm Raporu}
\fancyhead[R]{\small\color{slate400}Hafta 1: Marka \& Yerleşim Refaktörü}
\fancyfoot[C]{\small\color{slate400}\thepage}
\renewcommand{\headrulewidth}{0.4pt}
\renewcommand{\headrule}{\hbox to\headwidth{\color{codeframe}\leaders\hrule height \headrulewidth\hfill}}

% ================================================
% TITLE STYLE
% ================================================
\titleformat{\section}{\large\bfseries\color{indigo}}{}{0em}{}[\vspace{2pt}\color{codeframe}\hrule\vspace{4pt}]
\titleformat{\subsection}{\normalsize\bfseries\color{white!80!black}}{}{0em}{}

% ================================================
% CUSTOM BOXES
% ================================================
\tcbuselibrary{skins,breakable}
\newtcolorbox{changesummary}[1][]{
  enhanced,
  colback=slate800!60,
  colframe=indigo!50,
  fonttitle=\bfseries\small\color{white},
  title={#1},
  top=6pt, bottom=6pt, left=8pt, right=8pt,
  breakable
}
\newtcolorbox{codediff}{
  enhanced,
  colback=slate900,
  colframe=slate700,
  top=4pt, bottom=4pt, left=6pt, right=6pt,
  breakable
}
\newtcolorbox{keypoint}{
  enhanced,
  colback=indigo!10,
  colframe=indigo!60,
  fonttitle=\bfseries\small\color{indigo},
  title={Teknik Karar},
  top=6pt, bottom=6pt, left=8pt, right=8pt,
}

% ================================================
% DOCUMENT
% ================================================
\begin{document}

% ------------------------------------------------
% TITLE PAGE
% ------------------------------------------------
\begin{titlepage}
  \pagecolor{slate900}
  \color{white}
  \centering
  \vspace*{3cm}

  {\fontsize{10}{12}\selectfont\color{indigo}\bfseries\MakeUppercase{CodeAlchemist · SaaS Dönüşüm Serisi}}

  \vspace{1.5cm}
  {\fontsize{36}{40}\selectfont\bfseries Hafta 1}

  \vspace{0.5cm}
  {\fontsize{18}{22}\selectfont\color{slate400} Marka \& Yerleşim Refaktörü}

  \vspace{1cm}
  \color{codeframe}\hrule
  \vspace{1cm}

  \begin{tabular}{ll}
    \color{slate400}Tarih:    & \color{white} 6 Nisan 2026 \\[4pt]
    \color{slate400}Kapsam:   & \color{white} Frontend UI/UX · Design System · Gamification \\[4pt]
    \color{slate400}Durum:    & \color{indigo}\bfseries Tamamlandı ✓ \\
  \end{tabular}

  \vspace{2cm}
  \color{codeframe}
  \rule{\textwidth}{0.4pt}
  \vspace{0.5cm}

  {\small\color{slate400} Bu belge, CodeAlchemist'in \textit{öğrenci projesi} aşamasından çıkarılarak\\
  piyasaya sürülebilir bir SaaS ürününe dönüştürülmesi sürecinin\\
  \textbf{1. Hafta} kod tabanı değişikliklerini kapsamaktadır.}

  \vfill
  {\footnotesize\color{slate700} Tüm kaynak kodlar \texttt{client/src/} ve \texttt{server/} altında yönetilmektedir.}
\end{titlepage}
\pagecolor{white}
\color{black}

% ------------------------------------------------
% TABLE OF CONTENTS
% ------------------------------------------------
\tableofcontents
\newpage

% ================================================
% 1. GİRİŞ
% ================================================
\section{Giriş ve Motivasyon}

CodeAlchemist, yerleşik bir LLM yönlendirme motoru (\textit{Auto-Routing}), token bazlı ekonomisi ve çoklu model karşılaştırma özelliklerine sahip bir AI-powered geliştirici asistanıdır. Bu hafta gerçekleştirilen çalışmaların amacı, ürünün görsel dilini ve mimari dağılımını \textbf{profesyonel SaaS standartlarına} taşımaktır.

\begin{changesummary}[Hafta 1 Değişiklikleri — Özet]
\begin{itemize}[noitemsep]
  \item \textbf{Tasarım Sistemi:} \texttt{index.css} üzerinde \textit{Slate-900 Dark} default tema ve \textit{Indigo-500} marka rengi tanımlandı.
  \item \textbf{Landing Page:} Fayda temelli (\textit{benefit-driven}) Hero bölümü, rakip karşılaştırma tablosu ve custom SVG ikonlar eklendi.
  \item \textbf{Gamification:} Büyük modal kaldırıldı; artık yalnızca Profil sekmesinde yer alıyor.
  \item \textbf{Prompt Şablonları:} Çok renkli gradyanlar kaldırılıp Slate/Indigo sade tasarıma geçildi.
  \item \textbf{ChatInterface:} Tüm Fuchsia renk referansları Indigo ile değiştirildi.
\end{itemize}
\end{changesummary}

% ================================================
% 2. TASARIM SİSTEMİ
% ================================================
\section{Tasarım Sistemi: \texttt{index.css}}

\subsection{Motivasyon}

Uygulama önceden her bileşende farklı renk paleti ve font kullanıyordu. Bu durum görsel tutarsızlığa ve "hobi projesi" hissine yol açıyordu. Çözüm olarak merkezi bir CSS değişken sistemi (\textit{design tokens}) oluşturuldu.

\subsection{Yeni Tasarım Token'ları}

\textbf{Dosya:} \texttt{client/src/index.css} — Satır 560--583

\begin{lstlisting}[language=CSS, caption={Yeni SaaS Design Token Sistemi - index.css:560}]
/* Theme Variables - New SaaS Design System */
:root {
  /* Dark Theme Default (Slate-900) */
  --bg-base: #0F172A;      /* slate-900 */
  --bg-surface: #1E293B;   /* slate-800 */
  --bg-elevated: #334155;  /* slate-700 */

  --text-primary: #F1F5F9;   /* slate-100 */
  --text-secondary: #94A3B8; /* slate-400 */
  --text-muted: #64748B;    /* slate-500 */

  /* Marka Rengi: Indigo-500 */
  --accent: #6366F1;
  --accent-hover: #4F46E5;

  /* Imza Gradient - Hero ve badge'lerde kullanilir */
  --accent-gradient: linear-gradient(135deg, #6366F1, #8B5CF6);

  --border-color: #334155;
}
\end{lstlisting}

\begin{keypoint}
\textbf{Neden Indigo?} VS Code, GitHub, Linear ve Raycast gibi developer araçlarının büyük çoğunluğu mavi/indigo spektrumunda çalışır. Bu, kullanıcı beyninde otomatik bir "güvenilir araç" çağrışımı oluşturur. Fuchsia/pembe ise oyun ve tüketici uygulamalarıyla özdeşleştiğinden SaaS pazarı için uygun değildir.
\end{keypoint}

\subsection{Tipografi Standardizasyonu}

\textbf{Dosya:} \texttt{client/src/index.css} — Satır 8--14

\begin{lstlisting}[language=CSS, caption={Font Standardizasyonu - index.css:8}]
/* Uygulama geneli: Inter (Modern SaaS standardi) */
* {
  font-family: 'Inter', -apple-system, BlinkMacSystemFont,
               'Segoe UI', Roboto, sans-serif;
}

/* Kod bloklari: JetBrains Mono */
code, pre, .font-mono {
  font-family: 'JetBrains Mono', monospace !important;
}
\end{lstlisting}

% ================================================
% 3. LANDING PAGE
% ================================================
\section{Landing Page Refaktörü: \texttt{LandingPage.jsx}}

\subsection{Hero Bölümü ve "Radial Glow" Arka Planı}

\textbf{Dosya:} \texttt{client/src/components/LandingPage.jsx} — Satır 121--130

Arka plan katmanına yumuşak bir Indigo parlaması (radial gradient) eklenerek sayfanın derinliği artırıldı. Bu teknik, Linear, Vercel ve Supabase gibi premium SaaS ürünlerde yaygın kullanılır.

\begin{lstlisting}[language=JSX, caption={Radial Glow Arka Plan Efekti - LandingPage.jsx:121}]
{/* Background Blobs & Glows */}
<div className="fixed inset-0 overflow-hidden pointer-events-none">
  {/* YENI: Radial indigo glow - sayfa derinligini arttirir */}
  <div className="absolute inset-0
    bg-[radial-gradient(circle_at_50%_20%,
      rgba(99,102,241,0.25),
      transparent_60%)]" />

  {/* Ambient blob'lar */}
  <div className="absolute top-[-10%] left-[-10%] w-[40%] h-[40%]
       bg-indigo-600/10 blur-[120px] rounded-full animate-pulse" />
  <div className="absolute bottom-[-10%] right-[-10%] w-[40%] h-[40%]
       bg-purple-600/10 blur-[120px] rounded-full animate-pulse"
       style={{ animationDelay: '2s' }} />
</div>
\end{lstlisting}

\subsection{Fayda Temelli Hero Metni}

Eski "jenerik AI platform" mesajı yerine, geliştiricinin yaşadığı gerçek acı noktasına (\textit{pain point}) dokunan ve net bir çözüm sunan kelimeler yazıldı.

\begin{table}[h!]
\centering
\begin{tabular}{p{5.5cm}p{7cm}}
\toprule
\textbf{Eski Metin (Jenerik)} & \textbf{Yeni Metin (Fayda Odaklı)} \\
\midrule
"AI powered platform" & "5 Model aynı anda çalışsın." \\
"Cutting edge technology" & "En iyi cevabı sen seç." \\
"Innovative workflow" & "GPT-4o, Claude ve Gemini'yi tek soruda karşılaştır. Token bazlı fiyatlandırma ile sadece kullandığın kadar öde." \\
\bottomrule
\end{tabular}
\caption{Hero Mesaj Dönüşümü}
\end{table}

\textbf{Dosya:} \texttt{client/src/components/LandingPage.jsx} — Satır 157--176

\begin{lstlisting}[language=JSX, caption={Fayda Temelli Hero Baslik - LandingPage.jsx:157}]
<h1 className="text-5xl md:text-8xl font-black tracking-tight
               mb-6 leading-[0.85]">
  {/* Birinci satir: Beyaz - bold sorun tanimi */}
  <span className="block text-white">
    5 Model aynı anda çalışsın.
  </span>
  {/* Ikinci satir: Imza gradient - cozum vurgusu */}
  <span className="bg-clip-text text-transparent
                   bg-[var(--accent-gradient)]">
    En iyi cevabı sen seç.
  </span>
</h1>

{/* Destekleyici metin - spesifik, olcumlenebilir */}
<p className="text-base md:text-xl text-slate-400 max-w-2xl
              mx-auto mb-10 font-medium leading-relaxed">
  GPT-4o, Claude ve Gemini'yi tek soruda karşılaştır.
  Tek modele bağlı kalma. En iyi cevabı seç.
  Token bazlı fiyatlandırma ile sadece kullandığın kadar
  öde. Kuruş israf etme.
</p>
\end{lstlisting}

\subsection{Rakip Karşılaştırma Tablosu}

\textbf{Dosya:} \texttt{client/src/components/LandingPage.jsx} — Satır 3--48 (\texttt{ComparisonTable} bileşeni)

\begin{lstlisting}[language=JSX, caption={Rakip Karsilastirma Tablosu - LandingPage.jsx:3}]
const ComparisonTable = () => (
  <div className="mt-24 max-w-5xl mx-auto overflow-hidden
                  rounded-3xl border border-white/10
                  bg-white/[0.02] backdrop-blur-sm">
    <table className="w-full text-left border-collapse">
      <thead>
        <tr className="bg-white/[0.03]">
          <th>Ozellik</th>
          <th>GitHub Copilot</th>
          {/* CodeAlchemist sutunu indigo ile vurgulanir */}
          <th className="text-indigo-400 bg-indigo-500/5">
            CodeAlchemist
          </th>
        </tr>
      </thead>
      <tbody>
        {/* Model Secimi */}
        <tr><td>Model Secimi</td>
            <td>Tek Model (GPT-4o)</td>
            <td className="bg-indigo-500/5 font-bold text-white">
              Hibrit (5+ Model)
            </td></tr>
        {/* Fiyatlandirma */}
        <tr><td>Fiyatlandirma</td>
            <td>$10/ay (Sabit)</td>
            <td className="bg-indigo-500/5 font-bold text-white">
              Token Bazli (Kullandigin Kadar)
            </td></tr>
        {/* TRY / Yerel Odeme */}
        <tr><td>Yerel Odeme (TRY)</td>
            <td className="text-gray-500">Desteklenmiyor</td>
            <td className="bg-indigo-500/5 font-bold text-white">
              iyzico / TL Destegi
            </td></tr>
      </tbody>
    </table>
  </div>
);
\end{lstlisting}

\subsection{Custom SVG İkon Sistemi (Emoji Eliminasyonu)}

Önceki versiyonda emoji (🧠, ⚡, 📊) kullanılıyordu. Bu, ürünü görsel olarak "öğrenci projesi" kategorisine sokuyordu. Yerine platform kimliğine özel, ince çizgili (\texttt{strokeWidth=1.5}) SVG ikonlar tasarlandı.

\textbf{Dosya:} \texttt{client/src/components/LandingPage.jsx} — Satır 92--128

\begin{lstlisting}[language=JSX, caption={Custom SVG Ikon Sistemi - LandingPage.jsx:92}]
const features = [
  {
    title: "Zekayı Serbest Bırakın",
    desc: "GPT-4o, Claude 3.5 ve Gemini'yi ayni anda sorgulayin...",
    // YENI: Katmanli geometrik SVG - coklu model mimarisini temsil eder
    icon: (
      <svg width="32" height="32" viewBox="0 0 24 24"
           fill="none" xmlns="http://www.w3.org/2000/svg">
        {/* Ust katman - Tam opak */}
        <path d="M12 4L4 8L12 12L20 8L12 4Z"
              stroke="currentColor" strokeWidth="1.5"
              strokeLinecap="round" strokeLinejoin="round"/>
        {/* Orta katman - %50 opak */}
        <path d="M4 12L12 16L20 12"
              stroke="currentColor" strokeWidth="1.5"
              strokeLinecap="round" strokeLinejoin="round"
              opacity="0.5"/>
        {/* Alt katman - %20 opak (derinlik hissi) */}
        <path d="M4 16L12 20L20 16"
              stroke="currentColor" strokeWidth="1.5"
              strokeLinecap="round" strokeLinejoin="round"
              opacity="0.2"/>
      </svg>
    )
  },
  {
    title: "Seffaf Token Ekonomisi",
    icon: (
      <svg width="32" height="32" viewBox="0 0 24 24" fill="none">
        {/* Kesik cizgi cevre - "adil, seffaf" metaforu */}
        <circle cx="12" cy="12" r="9" stroke="currentColor"
                strokeWidth="1.5" strokeDasharray="4 4"/>
        {/* Saat ibre - "zaman = para" */}
        <path d="M12 7V12L15 15" stroke="currentColor"
              strokeWidth="1.5" strokeLinecap="round"/>
        <circle cx="12" cy="12" r="3" fill="currentColor"
                fillOpacity="0.2"/>
      </svg>
    )
  },
  {
    title: "Evrensel Entegrasyon",
    icon: (
      <svg width="32" height="32" viewBox="0 0 24 24" fill="none">
        {/* Yildirim - enerji ve hiz metaforu */}
        <path d="M13 2L3 14H12L11 22L21 10H12L13 2Z"
              stroke="currentColor" strokeWidth="1.5"
              strokeLinecap="round" strokeLinejoin="round"/>
        {/* Dolgu - imzanin arkasindan siran is isigi */}
        <path d="M13 2L3 14H12L11 22L21 10H12L13 2Z"
              fill="currentColor" fillOpacity="0.1"/>
      </svg>
    )
  }
];
\end{lstlisting}

\subsection{Feature Kartları: Derinlik ve Mikro-Animasyonlar}

\textbf{Dosya:} \texttt{client/src/components/LandingPage.jsx} — Satır 195--240

\begin{lstlisting}[language=JSX, caption={Premium Kart Tasarimi - LandingPage.jsx:195}]
<div
  className="group relative p-10
             rounded-[2.5rem]           /* Yumusak kesimler */
             bg-white/[0.01]            /* Cok ince cam efekti */
             border border-white/5
             hover:border-indigo-500/20 /* Hover: Indigo kenarligi */
             transition-all duration-500
             hover:-translate-y-2       /* Hover: Yukari yukselme */
             overflow-hidden"
>
  {/* HOVER BACKLIGHT - Arka isigindan siran enerji */}
  <div className="absolute -inset-px
    bg-gradient-to-br from-indigo-500/10 via-transparent
    to-purple-500/10
    opacity-0 group-hover:opacity-100
    transition-opacity duration-500" />

  {/* SVG Ikon Kabi - Glassmorphism */}
  <div className="relative w-16 h-16 mb-8 rounded-2xl
    bg-gradient-to-br from-indigo-500/10 to-purple-500/10
    border border-white/10 flex items-center justify-center
    text-3xl
    group-hover:scale-110 group-hover:rotate-3
    transition-transform duration-500">
    <div className="absolute inset-0 bg-indigo-500/5 blur-xl
                    rounded-full
                    group-hover:bg-indigo-500/20
                    transition-colors" />
    <span className="relative drop-shadow-md text-indigo-300">
      {f.icon}
    </span>
  </div>

  {/* ALT VURGU CIZGISI - Hover'da yatay genisleme */}
  <div className="absolute bottom-0 left-1/2 -translate-x-1/2
    w-0 h-1
    bg-gradient-to-r from-transparent via-indigo-500/40
    to-transparent
    group-hover:w-full
    transition-all duration-700" />
</div>
\end{lstlisting}

% ================================================
% 4. GAMİFİCATION REFAKTÖRİ
% ================================================
\section{Gamification: Hibrit Model}

\subsection{Sorun: Büyük Modal — SaaS Hissini Bozuyor}

Önceki versiyonda gamification (XP, seviye, rozetler), tam ekran bir modal olarak sunuluyordu. Bu yaklaşım şu sorunları yaratıyordu:

\begin{itemize}[noitemsep]
  \item Profesyonel SaaS araçlarında (VS Code, Cursor, Linear) böyle bir yapı bulunmuyor.
  \item Geliştirici iş akışını kesintiye uğratıyor.
  \item Header'daki LVL/XP rozeti dikkat dağıtıcıydı.
\end{itemize}

\subsection{Çözüm: Profil İçine Entegrasyon}

\textbf{Dosya:} \texttt{client/src/App.jsx} — Satır 2100--2114 (Silinen Kod)

\begin{lstlisting}[language=JSX, caption={Kaldirilan: Header LVL/XP Badge - App.jsx}]
{/* KALDIRILDI: Header'daki XP/Level rozeti */}
{/* BU KOD ARTIK YOK */}
{user && (
  <div className="flex items-center gap-2 bg-gray-800/40
                  border border-gray-700/50 rounded-full
                  px-3 py-1.5">
    <span className="text-[10px] font-black text-gray-500
                     uppercase tracking-tighter">LVL</span>
    <span className="text-xs font-bold text-indigo-400">
      {user.level || 1}
    </span>
    <div className="w-12 h-1 bg-gray-700 rounded-full">
      <div className="h-full bg-[var(--accent-gradient)]"
           style={{ width: `${(user.xp % 100) || 0}%` }} />
    </div>
  </div>
)}
\end{lstlisting}

\textbf{Dosya:} \texttt{client/src/components/UserProfileModal.jsx}

Gamification artık profil modalının "Başarılar" sekmesine taşındı. Kullanıcı ihtiyaç duyduğunda burger menüden erişebilir, chat akışı hiç bölünmez.

\begin{lstlisting}[language=JSX, caption={Profil Modalinda Sekmeli Yapi (Ozet) - UserProfileModal.jsx}]
const TABS = [
  { id: 'overview',      label: 'Genel Bakis', icon: '...' },
  { id: 'achievements',  label: 'Basarilar',   icon: '...' },
  { id: 'appearance',    label: 'Gorunum',     icon: '...' },
  { id: 'account',       label: 'Hesap',       icon: '...' },
];

{/* Basarilar sekmesi: Gamification burada */}
{activeTab === 'achievements' && (
  <GamificationPanel
    user={user}
    token={token}
    apiBase={apiBase}
  />
)}

{/* Gorunum sekmesi: ThemeStore burada */}
{activeTab === 'appearance' && (
  <ThemeStore
    currentTheme={currentTheme}
    onThemeChange={onThemeChange}
  />
)}
\end{lstlisting}

% ================================================
% 5. PROMPT ŞABLONLARI
% ================================================
\section{Prompt Şablonları: \texttt{PromptTemplates.jsx}}

\subsection{Sorun: Gökkuşağı Gradyanlar}

Her şablon butonu farklı bir renkle boyanmıştı (\texttt{from-blue-500 to-cyan-500}, \texttt{from-red-500 to-orange-500} vb.). Bu durum chat ekranını doğrudan görsel gürültüye dönüştürüyordu.

\textbf{Dosya:} \texttt{client/src/components/PromptTemplates.jsx}

\begin{lstlisting}[language=JSX, caption={Onceki Hali - Cok Renkli Gradyanlar (Kaldirildi)}]
// ESKI KOD - ARTIK KULLANILMIYOR
const promptTemplates = [
  { icon: '...',  label: 'Explain',  color: 'from-blue-500 to-cyan-500' },
  { icon: '...',  label: 'Debug',    color: 'from-red-500 to-orange-500' },
  { icon: '...',  label: 'Optimize', color: 'from-purple-500 to-pink-500' },
  { icon: '...',  label: 'Test',     color: 'from-green-500 to-emerald-500'},
  ...
];

// Rengarenk buton
<button className={`bg-gradient-to-r ${template.color} text-white ...`}>
\end{lstlisting}

\begin{lstlisting}[language=JSX, caption={Yeni Hali - Sade Slate/Indigo Tasarim - PromptTemplates.jsx:42}]
// YENI KOD - Birlesik, sakin tasarim
<button
  className="
    flex-1 min-w-[110px]
    flex items-center justify-center gap-2
    px-3 py-2.5 rounded-xl text-xs font-semibold

    /* Varsayilan: Koyu gri, ince kenarlik */
    bg-slate-800/40 border border-slate-700/50 text-slate-300

    /* Hover: Sadece Indigo vurgusu aktive olur */
    hover:bg-indigo-500/10
    hover:border-indigo-500/30
    hover:text-white

    /* Gecisin kendisi de animasyonlu */
    transition-all duration-300 ease-out backdrop-blur-sm
    group
  "
>
  {/* Ikon hover'da buyuyor ve belirginlesiyyor */}
  <span className="opacity-70
                   group-hover:opacity-100
                   group-hover:scale-110
                   transition-all duration-300">
    {template.icon}
  </span>
  <span>{template.label}</span>
</button>
\end{lstlisting}

% ================================================
% 6. CHAT INTERFACE - RENK NORMALİZASYONU
% ================================================
\section{ChatInterface Renk Normalizasyonu}

\textbf{Dosya:} \texttt{client/src/components/ChatInterface.jsx}

\subsection{Fuchsia → Indigo Migrasyonu}

Önceki versiyonda birçok interaktif element Fuchsia (\texttt{\#D946EF}) rengi kullanıyordu. Bu renk, marka kimliğiyle örtüşmüyordu ve diğer framework renkleriyle çakışıyordu.

\begin{table}[h!]
\centering
\begin{tabular}{lll}
\toprule
\textbf{Element} & \textbf{Eski Sınıf} & \textbf{Yeni Sınıf} \\
\midrule
Loading Spinner & \texttt{bg-fuchsia-500} & \texttt{bg-[var(--accent-gradient)]} \\
Magic Fix Butonu & \texttt{from-purple-500 to-pink-500} & \texttt{bg-indigo-600 hover:bg-indigo-500} \\
Toast başarı ikonu & \texttt{border-fuchsia-500/50} & \texttt{border-indigo-500/50} \\
Kopyala hover & \texttt{hover:text-fuchsia-300} & \texttt{hover:text-indigo-300} \\
Ses çalma aktif & \texttt{text-fuchsia-400} & \texttt{text-indigo-400} \\
\bottomrule
\end{tabular}
\caption{ChatInterface Renk Migrasyonu}
\end{table}

\begin{lstlisting}[language=JSX, caption={Loading Spinner - Oncesi ve Sonrasi - ChatInterface.jsx:1601}]
{/* ONCESI: Fuchsia */}
<div className="... bg-fuchsia-500/20 ...">
  <div className="... bg-gradient-to-r from-fuchsia-500 to-purple-500 ..."/>
</div>

{/* SONRASI: Marka degiskenleri */}
<div className="absolute inset-0 bg-indigo-500/20 rounded-full
                animate-ping" />
<div className="absolute inset-2 bg-[var(--accent-gradient)]
                rounded-full animate-spin
                shadow-[0_0_15px_rgba(99,102,241,0.5)]">
  <div className="absolute top-1 left-1 w-2 h-2
                  bg-white rounded-full opacity-60" />
</div>
\end{lstlisting}

% ================================================
% 7. UYGULAMA YAPISI (APP.JSX)
% ================================================
\section{Uygulama Ana Yapısı: \texttt{App.jsx}}

\subsection{Varsayılan Tema: Dark Mode}

\textbf{Dosya:} \texttt{client/src/App.jsx}

Kullanıcı herhangi bir ayar yapmadan uygulamayı açtığında artık doğrudan \textbf{Dark (Slate-900)} tema ile karşılanmaktadır. Bu karar şu nedenlere dayanmaktadır: Developer araçlarının \%90'ı dark mode kullanır (VS Code, GitHub, Cursor, Linear, Raycast).

\begin{lstlisting}[language=JSX, caption={Default Dark Theme - App.jsx}]
// Varsayilan tema karar mantigi
const [theme, setTheme] = useState(() => {
  // Once localStorage'a bakilir (kullanicinin tercihine saygi)
  const saved = localStorage.getItem('codebrain_theme');
  if (saved) return saved;
  // Default: dark (developer deneyimine gore dogal)
  return 'dark';
});
\end{lstlisting}

\subsection{"Magical Tools" Menüsü: Modal → Profil Yönlendirmesi}

\textbf{Dosya:} \texttt{client/src/App.jsx}

Eski versiyonda "Theme Store" ve "My Rank" menü öğeleri ayrı modalleri açıyordu. Bu, tek sayfalı bir uygulama için kötü bir UX pratiğiydi. Şimdi her ikisi de doğrudan ilgili profil sekmesini açmaktadır.

\begin{lstlisting}[language=JSX, caption={Magical Tools Menüsü - Modal Yerine Profil - App.jsx}]
{/* ONCESI - Ayri modal aciyordu */}
{ label: 'Theme Store', onClick: () => setShowThemeStore(true) }
{ label: 'My Rank',     onClick: () => setShowGamification(true) }

{/* SONRASI - Profil modalinin ilgili sekmesini aciyor */}
{
  label: 'Theme Store',
  onClick: () => {
    setProfileDefaultTab('appearance'); // 'appearance' sekmesi
    setShowProfile(true);               // Profili ac
  }
},
{
  label: 'My Rank',
  onClick: () => {
    setProfileDefaultTab('achievements'); // 'achievements' sekmesi
    setShowProfile(true);
  }
},
\end{lstlisting}

% ================================================
% 8. DEĞİŞİKLİK ÖZETİ TABLOSU
% ================================================
\section{Değişen Dosyalar — Özet Tablosu}

\begin{longtable}{p{5cm}p{2.5cm}p{6cm}}
\toprule
\textbf{Dosya} & \textbf{Değişim} & \textbf{Açıklama} \\
\midrule
\texttt{client/src/index.css} & Düzenlendi & Design token sistemi eklendi. Slate-900 default, Indigo-500 accent. \\
\texttt{client/src/App.jsx} & Düzenlendi & Dark mode default, Header LVL badge silindi, modal → profil yönlendirmesi. \\
\texttt{client/src/components/LandingPage.jsx} & Yeniden Yazıldı & Radial glow, benefit-driven hero, comparison table, custom SVG ikonlar, premium kart animasyonları. \\
\texttt{client/src/components/PromptTemplates.jsx} & Düzenlendi & Çok renkli gradyanlar kaldırıldı; Slate/Indigo sade tasarıma geçildi. \\
\texttt{client/src/components/ChatInterface.jsx} & Düzenlendi & Fuchsia → Indigo renk migrasyonu. Tüm marka renkleri CSS değişkenleriyle normalize edildi. \\
\texttt{client/src/components/UserProfileModal.jsx} & Genişletildi & Sekmeli yapı (Overview, Achievements, Appearance, Account). GamificationPanel ve ThemeStore entegre edildi. \\
\bottomrule
\end{longtable}

% ================================================
% 9. MİMARİ KARARLAR
% ================================================
\section{Hafta 1 Mimari Kararlar}

\begin{keypoint}
\textbf{"Kopyala-Yapıştır Renk" Yasağı:} Doğrudan Tailwind sınıf renkleri (\texttt{bg-indigo-600}) yerine mümkün olan her yerde CSS değişkenleri (\texttt{bg-[var(--accent)]}) kullanıldı. Bu, tema değişimini tek noktadan yönetmeyi mümkün kılıyor.
\end{keypoint}

\begin{keypoint}
\textbf{Gamification Görünürlük Prensibi:} "Motivasyon aracı dikkat dağıtmak için değil, bağlılığı artırmak için vardır." Kurala göre: Gamification \textit{mevcut} ama \textit{göze çarpmaz}. Kullanıcı istediğinde profil sekmesinden erişir.
\end{keypoint}

\begin{keypoint}
\textbf{SVG İkon Tasarım Dili:} \texttt{strokeWidth="1.5"} ve opacity katmanları (\texttt{opacity="0.5"}, \texttt{opacity="0.2"}) kullanılarak tek bir icon içinde derinlik hissi yaratıldı. Bu yaklaşım, Linear ve Figma'nın ikon sisteminden ilham almıştır.
\end{keypoint}

% ================================================
% 10. SONUÇ VE SONRAKİ HAFTA
% ================================================
\section{Sonuç ve Hafta 2 Planı}

Hafta 1 kapsamında CodeAlchemist'in görsel kimliği ve bileşen mimarisi profesyonel SaaS standartlarına kavuşturuldu. Ürün artık şu kriterleri karşılamaktadır:

\begin{itemize}[noitemsep]
  \item \textbf{Tanınabilir Bir Marka:} Slate-900 + Indigo-500 + JetBrains Mono kombinasyonu tutarlı bir developer kimliği oluşturuyor.
  \item \textbf{Net Değer Önerisi:} Landing page'de "neden Copilot değil?" sorusuna direkt cevap veriliyor.
  \item \textbf{Temiz UX:} Gereksiz modal katmanları temizlendi, her şey profil altında birleştirildi.
  \item \textbf{Ölçeklenebilir Design System:} Yeni özellikler, mevcut token sistemi üzerine kolayca inşa edilebilir.
\end{itemize}

\vspace{1cm}

\begin{changesummary}[Hafta 2 Planı: API Keys \& VS Code Eklentisi]
\begin{enumerate}[noitemsep]
  \item \textbf{Veritabanı:} \texttt{ApiKey} modeli ekleme (\texttt{server/models.py})
  \item \textbf{Backend:} \texttt{/v1/ask} endpoint'inin API Key auth ile stabilizasyonu (\texttt{server/app.py})
  \item \textbf{Profil UI:} "Developer Settings" sekmesi — key oluştur, isimlendir, iptal et
  \item \textbf{VS Code Extension:} \texttt{yo code} ile TypeScript iskeleti, Command Palette entegrasyonu
\end{enumerate}
\end{changesummary}

\vfill

\noindent\color{slate700}\rule{\textwidth}{0.4pt}
\vspace{4pt}
{\small\color{slate400}
\textit{CodeAlchemist SaaS Transformation Report} · Hafta 1/12 ·
Tüm kaynak kodlar \texttt{Dilakemer/code\_alchemist\_last\_version} deposunda bulunmaktadır.}

\end{document}
